\documentclass[11pt]{article}

\usepackage[a4paper,margin=1in]{geometry}
\usepackage[T1]{fontenc}
\usepackage{lmodern}
\usepackage{microtype}
\usepackage{amsmath,amssymb,amsthm,mathtools}

\newtheorem{theorem}{Theorem}
\newtheorem{remark}{Remark}

\newcommand{\X}{\mathcal X}
\newcommand{\A}{\mathcal A}
\newcommand{\Pcal}{\mathcal P}
\newcommand{\KL}{\operatorname{KL}}
\newcommand{\Breg}{\operatorname{B}}
\newcommand{\Dphi}{\operatorname{D}}
\newcommand{\dd}{\mathrm d}
\newcommand{\1}{\mathbf 1}

\title{Bregman Divergences, Csisz\'ar Divergences, and the Uniqueness of KL}
\author{}
\date{}

\begin{document}

\maketitle

\section{General measure-theoretic setting}

Let \((\X,\A)\) be a measurable space. We write \(\alpha\) and \(\beta\) for
probability measures on \((\X,\A)\).

For the Bregman part, it is convenient to work on a dominated model. Thus fix a
finite reference measure \(\lambda\) on \((\X,\A)\), and consider probability
measures of the form
\[
\alpha=p\lambda,
\qquad
\beta=q\lambda,
\]
where
\[
p=\frac{\dd\alpha}{\dd\lambda},
\qquad
q=\frac{\dd\beta}{\dd\lambda}.
\]
To avoid boundary issues, define the positive bounded part of the model by
\[
\Pcal_{\lambda,*}
=
\left\{
p\lambda:
p\in L^\infty(\lambda),\
\int_\X p\,\dd\lambda=1,\
\text{and there exist }0<m<M<\infty
\text{ such that }m\leq p\leq M
\right\}.
\]
The tangent directions are signed densities \(h\in L^\infty(\lambda)\) satisfying
\[
\int_\X h\,\dd\lambda=0.
\]
If \(p\in\Pcal_{\lambda,*}\) and \(h\) is such a tangent direction, then
\(p+\varepsilon h\) is again positive for all sufficiently small
\(\varepsilon\).

\section{Bregman divergence on a space of measures}

Let \(\Phi\) be a convex functional on a convex class of probability measures. If
\(\Phi\) is differentiable at \(\beta\) in the signed direction
\(\alpha-\beta\), its Bregman divergence is
\[
\Breg_\Phi(\alpha\mid\beta)
=
\Phi(\alpha)-\Phi(\beta)
-
D\Phi(\beta)[\alpha-\beta].
\]
In the dominated notation \(\alpha=p\lambda\), \(\beta=q\lambda\), this is
\[
\Breg_\Phi(p\lambda\mid q\lambda)
=
\Phi(p)-\Phi(q)-D\Phi(q)[p-q].
\]
Here
\[
D\Phi(q)[h]
=
\left.
\frac{\dd}{\dd\varepsilon}
\Phi(q+\varepsilon h)
\right|_{\varepsilon=0}
\]
is the first variation of \(\Phi\) at \(q\) in the direction \(h\).

By convexity,
\[
\Breg_\Phi(\alpha\mid\beta)\geq 0.
\]
Also, adding an affine functional to \(\Phi\) does not change
\(\Breg_\Phi\).

\section{Csisz\'ar, or \(f\)-divergence}

Let
\[
\varphi:(0,\infty)\to\mathbb R
\]
be convex and satisfy
\[
\varphi(1)=0.
\]
If \(\alpha\ll\beta\), the Csisz\'ar divergence generated by \(\varphi\) is
\[
\Dphi_\varphi(\alpha\mid\beta)
=
\int_\X
\varphi\!\left(
\frac{\dd\alpha}{\dd\beta}
\right)
\,\dd\beta.
\]
In dominated form, if \(\alpha=p\lambda\) and \(\beta=q\lambda\) with
\(q>0\), then
\[
\frac{\dd\alpha}{\dd\beta}
=
\frac{p}{q},
\]
and therefore
\[
\Dphi_\varphi(\alpha\mid\beta)
=
\int_\X
q(x)\,
\varphi\!\left(
\frac{p(x)}{q(x)}
\right)
\,\dd\lambda(x).
\]

For arbitrary probability measures \(\alpha,\beta\), one can use the usual
Lebesgue decomposition
\[
\alpha=\alpha_{\mathrm{ac}}+\alpha_{\mathrm{s}},
\qquad
\alpha_{\mathrm{ac}}\ll\beta,
\qquad
\alpha_{\mathrm{s}}\perp\beta.
\]
Writing
\[
r=\frac{\dd\alpha_{\mathrm{ac}}}{\dd\beta}
\]
and
\[
\varphi'_\infty
=
\lim_{t\to\infty}\frac{\varphi(t)}{t},
\]
the extended Csisz\'ar divergence is
\[
\Dphi_\varphi(\alpha\mid\beta)
=
\int_\X \varphi(r)\,\dd\beta
+
\varphi'_\infty\,\alpha_{\mathrm{s}}(\X).
\]
In the positive dominated model \(\Pcal_{\lambda,*}\), the singular term is absent.

\section{The KL example}

Define the entropy functional relative to \(\lambda\) by
\[
H_\lambda(p)
=
\int_\X p\log p\,\dd\lambda.
\]
Its first variation is
\[
DH_\lambda(q)[h]
=
\int_\X (1+\log q)h\,\dd\lambda.
\]
Therefore
\begin{align*}
\Breg_{H_\lambda}(p\lambda\mid q\lambda)
&=
\int_\X p\log p\,\dd\lambda
-
\int_\X q\log q\,\dd\lambda
-
\int_\X (1+\log q)(p-q)\,\dd\lambda \\
&=
\int_\X p\log p\,\dd\lambda
-
\int_\X p\log q\,\dd\lambda
-
\int_\X p\,\dd\lambda
+
\int_\X q\,\dd\lambda \\
&=
\int_\X p\log\frac{p}{q}\,\dd\lambda.
\end{align*}
Since \(p=\dd\alpha/\dd\lambda\) and \(q=\dd\beta/\dd\lambda\), this is
\[
\Breg_{H_\lambda}(\alpha\mid\beta)
=
\int_\X
\log\!\left(
\frac{\dd\alpha}{\dd\beta}
\right)
\,\dd\alpha
=
\KL(\alpha\mid\beta).
\]

On the other hand, if
\[
\varphi(t)=t\log t,
\]
then
\[
\Dphi_\varphi(\alpha\mid\beta)
=
\int_\X
q
\left(
\frac{p}{q}\log\frac{p}{q}
\right)
\,\dd\lambda
=
\int_\X p\log\frac{p}{q}\,\dd\lambda
=
\KL(\alpha\mid\beta).
\]
Thus KL is both a Bregman divergence and a Csisz\'ar divergence.

\section{Uniqueness theorem}

\begin{theorem}
Assume that \((\X,\A,\lambda)\) is nontrivial, in the sense that there is a
measurable set \(A\in\A\) with
\[
0<\lambda(A)<\lambda(\X)<\infty.
\]
Let \(\Phi\) be a convex functional on \(\Pcal_{\lambda,*}\), twice
Gateaux differentiable along bounded zero-mass directions. Let
\[
\varphi\in C^2(0,\infty)
\]
be convex and satisfy \(\varphi(1)=0\).

Suppose that
\[
\Breg_\Phi(\alpha\mid\beta)
=
\Dphi_\varphi(\alpha\mid\beta)
\]
for all \(\alpha,\beta\in\Pcal_{\lambda,*}\). Then there exist constants
\(c\geq 0\) and \(a\in\mathbb R\) such that
\[
\varphi(t)
=
c\,t\log t+a(t-1).
\]
Consequently,
\[
\Breg_\Phi(\alpha\mid\beta)
=
\Dphi_\varphi(\alpha\mid\beta)
=
c\,\KL(\alpha\mid\beta).
\]
Moreover,
\[
\Phi
=
c\,H_\lambda+\ell
\]
on \(\Pcal_{\lambda,*}\), where \(\ell\) is affine. Hence, apart from the
trivial zero divergence and scalar multiplication, the only common divergence is
KL.
\end{theorem}

\begin{proof}
Fix \(\beta=q\lambda\in\Pcal_{\lambda,*}\) and let
\(h\in L^\infty(\lambda)\) satisfy
\[
\int_\X h\,\dd\lambda=0.
\]
For sufficiently small \(\varepsilon\), define
\[
\alpha_\varepsilon=(q+\varepsilon h)\lambda.
\]
Then \(\alpha_\varepsilon\in\Pcal_{\lambda,*}\).

By the assumed equality of the two divergences,
\[
\Breg_\Phi(\alpha_\varepsilon\mid\beta)
=
\Dphi_\varphi(\alpha_\varepsilon\mid\beta).
\]

First compute the second variation of the Bregman side. Since
\[
\Breg_\Phi(\alpha_\varepsilon\mid\beta)
=
\Phi(q+\varepsilon h)-\Phi(q)-\varepsilon D\Phi(q)[h],
\]
we get
\[
\left.
\frac{\dd^2}{\dd\varepsilon^2}
\Breg_\Phi(\alpha_\varepsilon\mid\beta)
\right|_{\varepsilon=0}
=
D^2\Phi(q)[h,h].
\]

Now compute the second variation of the Csisz\'ar side:
\begin{align*}
\Dphi_\varphi(\alpha_\varepsilon\mid\beta)
&=
\int_\X
q\,
\varphi\!\left(
\frac{q+\varepsilon h}{q}
\right)
\,\dd\lambda \\
&=
\int_\X
q\,
\varphi\!\left(
1+\varepsilon\frac{h}{q}
\right)
\,\dd\lambda.
\end{align*}
Differentiating twice at \(\varepsilon=0\) gives
\[
\left.
\frac{\dd^2}{\dd\varepsilon^2}
\Dphi_\varphi(\alpha_\varepsilon\mid\beta)
\right|_{\varepsilon=0}
=
\varphi''(1)
\int_\X
\frac{h^2}{q}
\,\dd\lambda.
\]
Thus
\[
D^2\Phi(q)[h,h]
=
\varphi''(1)
\int_\X
\frac{h^2}{q}
\,\dd\lambda.
\]
Set
\[
c=\varphi''(1).
\]
Since \(\varphi\) is convex,
\[
c\geq 0.
\]

The entropy functional satisfies
\[
D^2H_\lambda(q)[h,h]
=
\int_\X \frac{h^2}{q}\,\dd\lambda.
\]
Therefore, for
\[
\Psi=\Phi-cH_\lambda,
\]
we have
\[
D^2\Psi(q)[h,h]=0
\]
for every \(q\in\Pcal_{\lambda,*}\) and every bounded zero-mass direction \(h\).

Now take arbitrary \(p\lambda,q\lambda\in\Pcal_{\lambda,*}\), and define the
line segment
\[
p_t=(1-t)q+tp,
\qquad
0\leq t\leq 1.
\]
Then \(p_t\lambda\in\Pcal_{\lambda,*}\), and
\[
\int_\X (p-q)\,\dd\lambda=0.
\]
Hence
\[
\frac{\dd^2}{\dd t^2}\Psi(p_t)
=
D^2\Psi(p_t)[p-q,p-q]
=
0.
\]
Thus \(t\mapsto \Psi(p_t)\) is affine. In particular,
\[
\Psi(p)-\Psi(q)-D\Psi(q)[p-q]=0.
\]
That is,
\[
\Breg_\Psi(p\lambda\mid q\lambda)=0.
\]
Since
\[
\Phi=cH_\lambda+\Psi,
\]
we obtain
\[
\Breg_\Phi(p\lambda\mid q\lambda)
=
c\,\Breg_{H_\lambda}(p\lambda\mid q\lambda)
=
c\,\KL(p\lambda\mid q\lambda).
\]

It remains to identify \(\varphi\). Since
\[
\Dphi_\varphi(\alpha\mid\beta)
=
c\,\KL(\alpha\mid\beta)
=
\Dphi_{c\,t\log t}(\alpha\mid\beta),
\]
we have
\[
\Dphi_g(\alpha\mid\beta)=0
\]
for all \(\alpha,\beta\in\Pcal_{\lambda,*}\), where
\[
g(t)=\varphi(t)-c\,t\log t.
\]
Also,
\[
g(1)=0.
\]

Choose arbitrary numbers
\[
0<x<1<y.
\]
Set
\[
\theta=\frac{y-1}{y-x}.
\]
Then
\[
0<\theta<1
\]
and
\[
\theta x+(1-\theta)y=1.
\]

Choose \(A\in\A\) with
\[
0<\lambda(A)<\lambda(\X),
\]
and put \(A^c=\X\setminus A\). Define
\[
q
=
\frac{\theta}{\lambda(A)}\1_A
+
\frac{1-\theta}{\lambda(A^c)}\1_{A^c},
\]
and
\[
p
=
x\,\frac{\theta}{\lambda(A)}\1_A
+
y\,\frac{1-\theta}{\lambda(A^c)}\1_{A^c}.
\]
Then
\[
\int_\X q\,\dd\lambda=1,
\]
and, because \(\theta x+(1-\theta)y=1\),
\[
\int_\X p\,\dd\lambda=1.
\]
Moreover,
\[
\frac{p}{q}
=
x
\quad\text{on }A,
\qquad
\frac{p}{q}
=
y
\quad\text{on }A^c.
\]
Since \(\Dphi_g(p\lambda\mid q\lambda)=0\), we get
\[
0
=
\theta g(x)+(1-\theta)g(y).
\]
Substituting
\[
\theta=\frac{y-1}{y-x},
\qquad
1-\theta=\frac{1-x}{y-x},
\]
we obtain
\[
0
=
\frac{y-1}{y-x}g(x)
+
\frac{1-x}{y-x}g(y).
\]
Multiplying by \(y-x\),
\[
(y-1)g(x)+(1-x)g(y)=0.
\]
Equivalently,
\[
(y-1)g(x)=(x-1)g(y),
\]
and hence
\[
\frac{g(x)}{x-1}
=
\frac{g(y)}{y-1}.
\]
Because \(0<x<1<y\) were arbitrary, the ratio
\[
\frac{g(t)}{t-1}
\]
is constant on \((0,\infty)\setminus\{1\}\). Thus there exists \(a\in\mathbb R\)
such that
\[
g(t)=a(t-1)
\]
for every \(t>0\). Therefore
\[
\varphi(t)
=
c\,t\log t+a(t-1).
\]
This proves the claim.
\end{proof}

\section{Normalizations}

The generator of a Csisz\'ar divergence is not unique: adding a linear term
\(a(t-1)\) does not change the divergence, because
\[
\int_\X
q
\left(
\frac{p}{q}-1
\right)
\,\dd\lambda
=
\int_\X p\,\dd\lambda-\int_\X q\,\dd\lambda
=
0.
\]
Thus
\[
\varphi(t)=c\,t\log t+a(t-1)
\]
and
\[
\varphi(t)=c\,t\log t
\]
define the same divergence on probability measures.

If one imposes the common normalization
\[
\varphi'(1)=0,
\]
then
\[
0=\varphi'(1)=c+a,
\]
so
\[
a=-c,
\]
and
\[
\varphi(t)
=
c\bigl(t\log t-t+1\bigr).
\]
If one also imposes
\[
\varphi''(1)=1,
\]
then \(c=1\), and the normalized generator is
\[
\varphi(t)=t\log t-t+1.
\]
It gives exactly
\[
\Dphi_\varphi(\alpha\mid\beta)
=
\KL(\alpha\mid\beta).
\]

\section{Independence of the reference measure}

Although \(H_\lambda\) uses a reference measure, its Bregman divergence does not.
Indeed, suppose \(\lambda'\) is another equivalent reference measure, with
\[
\dd\lambda'=s\,\dd\lambda,
\qquad
s>0.
\]
If \(\mu=p\lambda=p'\lambda'\), then
\[
p'=\frac{p}{s}.
\]
Hence
\begin{align*}
H_{\lambda'}(\mu)
&=
\int_\X p'\log p'\,\dd\lambda' \\
&=
\int_\X \frac{p}{s}\log\!\left(\frac{p}{s}\right)s\,\dd\lambda \\
&=
\int_\X p\log p\,\dd\lambda
-
\int_\X p\log s\,\dd\lambda.
\end{align*}
Thus
\[
H_{\lambda'}(\mu)
=
H_\lambda(\mu)
-
\int_\X \log s\,\dd\mu.
\]
The difference between \(H_{\lambda'}\) and \(H_\lambda\) is affine in \(\mu\),
so their Bregman divergences are identical. Therefore
\[
\Breg_{H_\lambda}(\alpha\mid\beta)
=
\KL(\alpha\mid\beta)
\]
is a measure-theoretic object, not an artifact of the choice of \(\lambda\).

\begin{remark}
For arbitrary probability measures, the KL divergence is understood with the usual
extended convention
\[
\KL(\alpha\mid\beta)
=
\begin{cases}
\displaystyle
\int_\X
\log\!\left(\frac{\dd\alpha}{\dd\beta}\right)
\,\dd\alpha,
& \alpha\ll\beta, \\[1.2em]
+\infty,
& \alpha\not\ll\beta.
\end{cases}
\]
For the nonzero common divergence, the corresponding Csisz\'ar generator is
\[
\varphi(t)=c\,t\log t+a(t-1),
\qquad c>0.
\]
Its extended Csisz\'ar divergence is
\[
\Dphi_\varphi(\alpha\mid\beta)
=
c\,\KL(\alpha\mid\beta).
\]
\end{remark}

\end{document}